where @@M0@@ is the ring @@M1@@, filtered by weight with respect to the @@M2@@, @@M3@@ having weight @@M4@@, and for the element @@M5@@.

We embed @@M6@@ in @@M7@@ by the symmetric functions, and thus obtain an injective homomorphism: @@M8@@; we then have:
@@M9@@
We deduce, by (6.2.1):
@@M10@@
and obtain:
@@M11@@
Finally,
@@M12@@
Now, by (6.5.3), @@M13@@
that is: @@M14@@.

\sgaStableTarget{sga6.e05.cor.6-6-2}\textbf{Corollary 6.6.2.} \uline{Relation} 6.6.1.iii) \uline{holds for every element} @@M15@@ \uline{of augmentation} @@M16@@.

To prove this, we shall use the lemma:

\sgaStableTarget{sga6.e05.lem.6-6-3}\textbf{Lemma 6.6.3.} \uline{For every} @@M17@@, @@M18@@ \uline{is a} @@M19@@-\uline{ideal of} @@M20@@.

We may suppose that @@M21@@. Let @@M22@@; we can write @@M23@@ where the @@M24@@ are the isobaric components of weight @@M25@@ of the @@M26@@ (3.6). For every @@M27@@, the polynomials @@M28@@ are therefore polynomials in the first @@M29@@ coefficients of @@M30@@. If @@M31@@, they are therefore zero, and consequently the @@M32@@ belong to @@M33@@, and
